\section{Discussion}\label{sec:discussion}

\paragraph{What the four studies show together.}
The composition identity says when a joint test can reject an alignment that \emph{every} single query accepts. It is
silent on single tests that draw only some of the single queries. In the small networks of Studies~1--3 that
distinction did little: joint-only failures were rare, and a joint test spent its budget less well than a single test
drawn from the full-domain panel. In GPT-2 it decided the result. The full-domain panel holds the queries in which the
swapped variable changes, which are also the queries DAS is trained on, and the position subspace that DAS finds in
GPT-2 also moves name information. A full-domain position query cannot show this, because it changes the position at
the same time; a position query whose donor keeps the position can, and so can every joint query, because its second
donor's names arrive through the position subspace too. Most of the joint tests' advantage in Study~4 therefore comes
from single queries that the prespecified comparison left out. Against single tests that include them the advantage
shrinks to \ContrastAllS{} in GPT-2 Small at \Budget{} calls, and it stays positive. Most joint-panel failures in GPT-2
are irreducible (\IrrShareS\% in GPT-2 Small and \IrrShareM\% in GPT-2 Medium, against \TrnIrrShare\% in Study~2),
yet no GPT-2 Small alignment passes every single query, and rank-2 and rank-4 DAS, which fail none of the full-domain
queries, fail the single position queries whose donor keeps the position.

\paragraph{What to test.}
For alignments of the kind studied here, six practical points follow.
\begin{enumerate}
  \item \emph{Test the input range the claim covers.} Alignments fitted on part of the range can pass every test
  drawn from that part and fail elsewhere, as all restricted lookup alignments of Study~1 and most restricted DAS
  alignments of Study~2 did. Where a few values already determine a variable's code, as in the additive programs,
  restricted fitting does no harm, but a test confined to that range cannot show this.
  \item \emph{Use donors that differ in more than the swapped variable.} A minimal pair cannot reveal an edit that
  disturbs the other variable; in Study~2 minimal-pair tests detected fewer failing alignments than full-domain
  single tests at each budget in \cref{fig:trained}b.
  \item \emph{Include single queries in which the swapped variable keeps its value.} These queries check that an
  intervention leaves the other variable alone, the isolation side of RAVEL's scores \citep{huang2024ravel}. In
  Study~4 a test drawn only from the queries DAS is trained on accepted alignments that fail \RateUnchangedAS\% of the
  position queries outside them. Joint queries are a cheap way to draw such queries, because every one of them changes
  both variables, and at small budgets they detected the most failing alignments in GPT-2.
  \item \emph{For orthogonal alignments at one site, use margins to find what joint tests can add beyond a complete
  single test.} Where single edits land on their targets, joint tests are redundant (\cref{prop:redundant}), and
  this can be checked. Approximate alignments essentially never land exactly; for them a joint query can fail only
  where the other edit's miss is large compared with the margins of both single queries it reduces to
  (\cref{prop:margin}). In Study~2 such failures were rare and sparse, and budget spent on joint queries detected
  fewer failing alignments than the same budget spent on full-domain single queries (\TrnContrastOneTwenty{}
  \TrnContrastOneTwentyCI{} at 120 calls). The signed first-order rule, computed from single queries and gradients
  alone, flagged \TrnMbRuleFlagged{} of Study~2's \TrnPerfectJointFailWord{} joint-only alignments and, fixed in advance
  and applied unchanged, all \MsPerfectJointFailWord{} of Study~3's (both computed post hoc); in Study~4 it flagged all
  \PerfectWordM{} of GPT-2 Medium's and \DFlagged{} of the \DJointOnly{} joint-only alignments of arm d, with no false
  flag. These counts are small, and the rule is meant only for alignments that pass their single queries, which in
  GPT-2 are few. Because one gradient per single query serves all joint queries built on it, the rule is cheap once a
  single census has been taken; at budgets of a few dozen queries it is not, and a joint test guided by it did worse
  than a plain single test (\cref{sec:strategies}).
  \item \emph{Where the variables sit at different sites, the reduction holds only through the upstream variable.} In
  Study~3, joint tests then detected failing alignments slightly more often than alternating single tests
  (\MsContrast{} \MsContrastCI{} at 120 calls), mostly, a post hoc analysis shows, because every joint query exercises
  the downstream variable, the weaker one there; against single tests aimed at that variable the difference is
  \MsMechWorse. MIB's alignment of the two IOI variables to whole attention heads of GPT-2 Small places them at
  different sites, two of the position heads a layer below the token head \citep{mueller2025mib}, which is closer to
  Study~3's setting. Under the counterfactual model of Study~4, however, the token head fails \HeadsRateB\% of its
  token queries, so single and joint tests both detect every such alignment, and the comparison between them is
  moot. This is a statement about that model of the task; MIB scores its IOI alignments differently.
  \item \emph{Do not use distance from natural states as a validity test.} Edits that move directions the network
  ignores are harmless, and raw distance rejected every valid alignment that made them.
\end{enumerate}

\paragraph{Overlapping subspaces.}
\Cref{lem:composition} needs orthogonal projectors. Subspaces fitted one variable at a time need not be orthogonal,
and then a joint intervention is not a composition of single ones and can expose interactions that single tests miss;
the general theory therefore requires joint interventions \citep{geiger2025abstraction}. Study~4 measured this in two
places. In Study~2's networks, separately fitted subspaces made joint-only failures three times as common as
Study~2's orthogonal alignments did (\DShare\% against \DRefShare\%), and the overlap separated the joint-only ones
(area under the ROC curve \DAuroc). In GPT-2 the separately fitted subspaces barely overlapped, and the joint failure
rate even fell with the overlap (Spearman \SepSpearman), so in that model the leak of names into the position subspace,
not the overlap, is what joint tests find. Invertible non-linear alignment maps that give each variable its own block
of latent coordinates, such as those of \citet{sutter2025nonlinear}, satisfy the identity of \cref{lem:composition} in
those coordinates, so \cref{prop:redundant,cor:irreducible} carry over to them; their difficulty is a different one,
since maps that flexible can align a network with almost any algorithm.

\paragraph{Limits.}
\Cref{lem:composition} needs orthogonal projectors at a single site. For variables at different sites it holds in
one direction only (\cref{rem:sites}), and Study~3 tests one pair of sites per architecture. The tasks of Studies~1--3
have two variables with three values each and 27 inputs, which is what makes exhaustive censuses possible. With so few
natural states an approximate edit essentially never lands on its target, so the redundancy check is informative only
for exact alignments; in larger domains tests are sampled, and whether single edits land has to be estimated. The
networks of Studies~1--3 are small: compiled networks give each variable its own coordinates by construction, and the
trained ones learn from 27 inputs. In Study~2 the full-domain DAS alignments are scored on the queries they were
trained on, so passing every single query measures their fit to those queries; refitted on half of them, they failed
the held-out half slightly more often, which moved the comparison slightly further from joint tests (post hoc). In
Study~3 the alignments at the MLP's second hidden layer were poor, so its estimate of the joint-only share rests on
\MsPerfect{} alignments, all on additive tasks. Study~4 has one task, with two and three values per variable and
\NTemplates{} templates per task. GPT-2 Small answered all \NPrompts{} prompts of only \GatePassedSmall{} of
\GateDrawnSmall{} held-out tasks with the required margin, so its confirmation rests on the tasks it finds easiest, and
on \GatePassedSmall{} rather than \NTasksConf{} tasks; GPT-2 Medium's rests on \NTasksMedium{}, set by the cost rule,
and its alignments were trained with TF32 matrix products. Larger models and other tasks may behave differently. The
size of the comparison at equal budgets depends on how queries are costed, and it answers which test to run on a fixed
budget; the add-on comparison answers what a joint test adds to a single one. Study~1's labels are exact for one
precise claim, and Studies~2--4 have no labels independent of the tests; none of the studies says which explanation
of a network is the true one \citep{meloux2025identifiable,chen2026rome}. Finally, the census of Study~1, the analyses
marked post hoc in Studies~2 and~3, including the margin analysis and the held-out split, the test strategies of
\cref{sec:strategies}, and Study~4's comparison with all single queries and its verifier were added after the frozen
results or the development results were known; they are computations on the saved networks and records and do not
change any frozen decision (\cref{app:gpt2-verification}). Study~3's predictions were fixed before its networks
existed and Study~4's before any model was run, but both studies were designed after the results of the studies before
them.
